% ============================================================================
%  VSEPR-SIM — §R: 32-Bit Identity Word, Hourglass Model, Lookglass Model
%  Printer-ready LaTeX.  Compile with pdflatex (twice for refs).
% ============================================================================
\documentclass[11pt, a4paper, oneside]{article}

% ── Geometry ────────────────────────────────────────────────────────────────
\usepackage[
  top=1.0in, bottom=1.0in, left=1.15in, right=1.15in,
  headheight=14pt
]{geometry}

% ── Fonts / encoding ───────────────────────────────────────────────────────
\usepackage[T1]{fontenc}
\usepackage{lmodern}
\usepackage{microtype}

% ── Math ───────────────────────────────────────────────────────────────────
\usepackage{amsmath, amssymb, amsthm, mathtools}
\usepackage{bm}

% ── Tables ─────────────────────────────────────────────────────────────────
\usepackage{booktabs}
\usepackage{array}
\usepackage{tabularx}

% ── Lists ──────────────────────────────────────────────────────────────────
\usepackage{enumitem}
\setlist{nosep, leftmargin=1.5em}

% ── Code listings ─────────────────────────────────────────────────────────
\usepackage{listings}
\usepackage{xcolor}

\definecolor{lstbg}{RGB}{248,248,248}
\definecolor{lstkw}{RGB}{0,0,180}
\definecolor{lstcm}{RGB}{80,120,80}
\definecolor{lstst}{RGB}{160,30,30}

\lstdefinestyle{cppstyle}{
  language=C++,
  basicstyle=\ttfamily\small,
  backgroundcolor=\color{lstbg},
  keywordstyle=\color{lstkw}\bfseries,
  commentstyle=\color{lstcm}\itshape,
  stringstyle=\color{lstst},
  frame=single, framesep=4pt,
  breaklines=true, keepspaces=true,
  tabsize=4,
  showstringspaces=false,
  numbers=left, numberstyle=\tiny\color{gray}, numbersep=6pt
}
\lstset{style=cppstyle}

% ── Theorem environments ───────────────────────────────────────────────────
\theoremstyle{definition}
\newtheorem{definition}{Definition}[section]
\newtheorem{postulate}[definition]{Postulate}
\newtheorem{remark}[definition]{Remark}

% ── Headers / footers ─────────────────────────────────────────────────────
\usepackage{fancyhdr}
\pagestyle{fancy}
\fancyhf{}
\fancyhead[L]{\small\textit{VSEPR-SIM Research Platform}}
\fancyhead[R]{\small\textit{32-Bit Model, Hourglass \& Lookglass}}
\fancyfoot[C]{\thepage}
\renewcommand{\headrulewidth}{0.4pt}

% ── Section formatting ─────────────────────────────────────────────────────
\usepackage{titlesec}
\titleformat{\section}
  {\Large\bfseries}{§\thesection}{1em}{}
\titleformat{\subsection}
  {\large\bfseries}{\thesubsection}{1em}{}
\titleformat{\subsubsection}
  {\normalsize\bfseries}{\thesubsubsection}{1em}{}

% ── Colours / boxes ────────────────────────────────────────────────────────
\usepackage{tcolorbox}
\tcbuselibrary{skins, breakable}
\definecolor{bitblue}{RGB}{220,235,250}
\definecolor{hgolden}{RGB}{255,245,220}
\definecolor{lgmint}{RGB}{220,250,235}

% ── Hyperlinks ─────────────────────────────────────────────────────────────
\usepackage{hyperref}
\hypersetup{
  colorlinks=true,
  linkcolor=blue!60!black,
  citecolor=green!50!black,
  urlcolor=blue!70!black
}

% ── Macros ─────────────────────────────────────────────────────────────────
\newcommand{\R}{\mathbb{R}}
\newcommand{\Z}{\mathbb{Z}}
\newcommand{\N}{\mathbb{N}}
\newcommand{\boldI}{\bm{I}}
\newcommand{\bits}[1]{\texttt{[#1~bits]}}
\newcommand{\file}[1]{\texttt{#1}}

% ============================================================================
\begin{document}
% ============================================================================

\begin{titlepage}
  \centering
  \vspace*{2cm}
  {\Huge\bfseries VSEPR-SIM Research Platform\par}
  \vspace{1cm}
  {\Large\bfseries §R\quad Three Foundational Models\par}
  \vspace{0.5cm}
  \rule{0.6\textwidth}{0.4pt}\par
  \vspace{0.8cm}
  {\large
    \textbf{I.}  The 32-Bit Particle Identity Word\par
    \vspace{0.3cm}
    \textbf{II.} The Hourglass Convergence Model\par
    \vspace{0.3cm}
    \textbf{III.}The Lookglass Bidirectional Feedback Model\par
  }
  \vspace{1.5cm}
  {\normalsize
    All three models operate on the Identity–State Decomposition
    framework introduced in §0.\par
    \vspace{0.5cm}
    Implementation files:\par
    \file{coarse\_grain/core/bit32\_identity.hpp}\par
    \file{coarse\_grain/models/hourglass\_model.hpp}\par
    \file{coarse\_grain/models/lookglass\_model.hpp}\par
  }
  \vfill
  {\small Status: Research-oriented active development}
\end{titlepage}

\tableofcontents
\clearpage

% ============================================================================
\section{The 32-Bit Particle Identity Word}
\label{sec:bit32}
% ============================================================================

\subsection{Motivation}

The §0 Identity–State Decomposition Framework defines a six-component
particle identity vector
%
\begin{equation}
  \boldI_i =
  \begin{bmatrix}
    Z_i \\ A_i \\ Q_i \\ \Sigma_i \\ \Lambda_i \\ \Theta_i
  \end{bmatrix}
  \label{eq:identity-vector}
\end{equation}
%
where $Z_i$ is atomic number, $A_i$ is mass identity, $Q_i$ is effective
charge participation, $\Sigma_i$ is structural role, $\Lambda_i$ is
stability class, and $\Theta_i$ is provenance memory.

In production simulations the identity vector is evaluated for every
bead at every step.  Storing it as six separate fields costs memory
bandwidth and breaks cache locality.  The 32-bit identity word packs
all six fields into a single \texttt{uint32\_t} that fits in a
register, enabling branch-free scalar read-outs and SIMD-friendly
bulk operations.

\subsection{Bit-Field Layout}

\begin{tcolorbox}[
  colback=bitblue, colframe=blue!40!black,
  title={32-Bit Identity Word Layout},
  fonttitle=\bfseries\small
]
\begin{center}
\ttfamily\small
\begin{tabular}{|c|c|c|c|c|c|}
\hline
\textbf{31--24} & \textbf{23--16} & \textbf{15--8} &
\textbf{7--5} & \textbf{4--3} & \textbf{2--0} \\
\hline
$Z_i$~[8] & $A_i$~[8] & $Q_i$~[8] & $\Sigma_i$~[3] & $\Lambda_i$~[2] & $\Theta_i$~[3] \\
\hline
\end{tabular}
\end{center}
\end{tcolorbox}

\noindent Total width: 8 + 8 + 8 + 3 + 2 + 3 = \textbf{32 bits exactly}.

\subsection{Field Encodings}
\label{sec:bit32-encodings}

\begin{center}
\begin{tabular}{llll}
\toprule
Field & Bits & Range & Encoding \\
\midrule
$Z_i$ & 8 & 0--118 & Atomic number; 0 = unassigned \\
$A_i$ & 8 & 0--255 & Mass bucket; 0 = natural-abundance average \\
$Q_i$ & 8 & $-32.0$\ldots$+31.75$\,e & Signed s8.2 fixed-point, resolution 0.25\,e \\
$\Sigma_i$ & 3 & 0--4 & StructuralRole enum (5 values; 3 bits sufficient) \\
$\Lambda_i$ & 2 & 0--3 & StabilityClass enum (4 values; 2 bits exact) \\
$\Theta_i$ & 3 & 0--7 & Provenance tag (8 provenance buckets) \\
\bottomrule
\end{tabular}
\end{center}

\subsubsection{Charge Encoding ($Q_i$)}

$Q_i$ uses signed 8-bit fixed-point with two fractional bits (s8.2),
giving a resolution of $0.25\,e$ and a range of $[-32.0, +31.75]\,e$.
This is sufficient for all common ions (Th$^{4+}$, Fe$^{3+}$,
organic anions) and avoids floating-point operations in the hot path.

The round-trip equation is:
%
\begin{equation}
  q_{\rm fp} = q_{\rm raw} \times 0.25\,e,
  \qquad
  q_{\rm raw} = \operatorname{round}(q_{\rm fp} / 0.25\,e)
  \label{eq:charge-roundtrip}
\end{equation}

\subsubsection{Provenance Tag ($\Theta_i$)}

The full provenance hash lives in \texttt{Bead::parent\_atom\_indices}.
The 3-bit tag is a coarse provenance bucket sufficient for routing
decisions without dereferencing the parent record:

\begin{center}
\begin{tabular}{cl}
\toprule
Tag & Meaning \\
\midrule
0 & Virgin / untracked \\
1 & Mapped from atomistic parent \\
2 & FIRE-relaxed minimum \\
3 & CG-only (no atomistic parent) \\
4 & Reaction product \\
5 & Crystal seed \\
6 & Foreign import \\
7 & Hypothesis / mutated structure \\
\bottomrule
\end{tabular}
\end{center}

\subsection{Round-Trip Guarantee}

\begin{definition}[Round-Trip Consistency]
  A 32-bit identity word $w$ is \emph{round-trip consistent} if
  $\operatorname{pack}(\operatorname{unpack}(w)) = w$.
\end{definition}

All valid inputs satisfy this property.  The only invalid words are
those with $\Sigma_i \in \{5,6,7\}$ (reserved), which are never
produced by \texttt{pack\_identity()}.

\subsection{Implementation}

The full implementation is in
\file{coarse\_grain/core/bit32\_identity.hpp}.
Key interface:

\begin{lstlisting}
// Pack from individual fields
Identity32 id = pack_identity(Z, A, Q_fp, sigma, lambda, theta);

// Pack directly from a Bead
Identity32 id = pack_from_bead(bead, Z);

// Read individual fields (branch-free bit shifts)
uint8_t        z   = id.Z();
double         q   = id.Q();        // decoded to floating-point
StructuralRole sig = id.sigma();
StabilityClass lam = id.lambda();
ProvenanceTag  th  = id.theta();

// Debug dump
std::cout << id.to_string() << '\n';
// → Identity32{ Z=6 A=0 Q=0 Σ=Directional-Covalent Λ=Ambient-Stable Θ=Mapped-Atomistic }

// Verify round-trip
assert(identity32_roundtrip_ok(id));
\end{lstlisting}

% ============================================================================
\section{The Hourglass Convergence Model}
\label{sec:hourglass}
% ============================================================================

\subsection{Concept}

In standard molecular optimisation, a single trajectory descends a
potential energy surface from an initial configuration to the nearest
local minimum.  This is efficient but blind: it finds \emph{a} minimum,
not necessarily the \emph{best} minimum accessible from the current state.

The Hourglass Model addresses this by running a population of candidate
configurations through a three-layer filter at each macro-step.  The
geometry of an hourglass is the natural metaphor:

\begin{center}
\begin{tabular}{cl}
\toprule
Layer & Role \\
\midrule
Wide mouth (H1) & Generate $N_{\rm cand}$ perturbed candidates \\
Narrow neck (H2) & Filter by stability, energy, and role weights \\
Wide base (H3) & Score and rank the survivors \\
\bottomrule
\end{tabular}
\end{center}

\subsection{Formal Specification}
\label{sec:hourglass-formal}

\subsubsection{H1 — Mouth: Candidate Generation}

Let $\mathbf{r}^{(0)}$ be the current bead configuration.
Generate $N_{\rm cand}$ candidates:
%
\begin{equation}
  \mathbf{r}^{(c)} = \mathbf{r}^{(0)} + \delta^{(c)},
  \quad \delta^{(c)}_i \sim \mathcal{U}(-\delta_{\max}, +\delta_{\max})^3,
  \quad c = 1,\ldots,N_{\rm cand}
  \label{eq:hg-mouth}
\end{equation}
%
Uniform perturbations (not Gaussian) ensure the sampled volume in
configuration space has a well-defined boundary at radius $\delta_{\max}$.

\subsubsection{H2 — Neck: Three-Gate Filter}

Each candidate passes three gates in sequence.  Rejection at any gate
eliminates the candidate immediately.

\paragraph{Gate 1 — Stability filter.}
%
\begin{equation}
  \bar{\Lambda}^{(c)} \;=\;
  \frac{1}{N} \sum_{i=1}^N \Lambda_i^{(c)}
  \;\geq\; \Lambda_{\min}
  \label{eq:gate1}
\end{equation}

\paragraph{Gate 2 — Energy filter.}
%
\begin{equation}
  \Delta E^{(c)} \;=\; E^{(c)} - E_{\rm basin} \;\leq\; E_{\rm tol},
  \qquad
  E_{\rm basin} = \min_c E^{(c)}
  \label{eq:gate2}
\end{equation}

\paragraph{Gate 3 — Role weight filter.}

Let $w_{\rm dom}^{(i)}$ denote the dominant channel weight of bead $i$:
%
\begin{equation}
  w_{\rm dom}^{(i)} = \max\!\bigl(w_{\rm steric}^{(i)},\;
                                   w_{\rm elec}^{(i)},\;
                                   w_{\rm disp}^{(i)}\bigr)
\end{equation}
%
Then:
%
\begin{equation}
  \frac{1}{N}\sum_{i=1}^N w_{\rm dom}^{(i)} \;\geq\; w_{\rm min}
  \label{eq:gate3}
\end{equation}

\subsubsection{H3 — Base: Composite Score}

Surviving candidates are ranked by:
%
\begin{equation}
  \mathcal{S}^{(c)} =
    w_E \cdot \frac{E_{\rm ref} - E^{(c)}}{|E_{\rm ref}| + \epsilon}
  + w_\Lambda \cdot \frac{\bar\Lambda^{(c)}}{3}
  + w_\Sigma \cdot \frac{1}{N}\sum_i w_{\rm dom}^{(i)}
  \label{eq:hg-score}
\end{equation}
%
The top $k \leq N_{\rm out}$ candidates are returned in descending score.

\subsection{Parameters}

\begin{center}
\begin{tabular}{llrl}
\toprule
Parameter & Symbol & Default & Meaning \\
\midrule
$N_{\rm cand}$ & — & 64 & Candidate count \\
$\delta_{\max}$ & — & 0.5\,\AA & Perturbation radius \\
$\Lambda_{\min}$ & — & 1 (Metastable) & Stability floor \\
$E_{\rm tol}$ & — & 5.0\,kcal/mol & Energy ceiling above basin \\
$w_{\rm min}$ & — & 0.5 & Role-weight floor \\
$N_{\rm out}$ & — & 8 & Maximum ranked survivors \\
$w_E, w_\Lambda, w_\Sigma$ & — & 1.0, 0.5, 0.3 & Score weights \\
\bottomrule
\end{tabular}
\end{center}

\subsection{Implementation}

The full implementation is in
\file{coarse\_grain/models/hourglass\_model.hpp}.

\begin{lstlisting}
HourglassParams params;
params.N_cand    = 128;
params.delta_max = 0.3;    // tighter sampling
params.E_tol     = 3.0;    // stricter energy gate

HourglassResult result = run_hourglass(sys, params);

if (result.any_survived) {
    const CandidateRecord* best = hourglass_best(result);
    // Apply best configuration to system...
}

// Full audit: all candidates and gate verdicts
for (const auto& c : result.candidates) {
    std::cout << "Candidate " << c.candidate_id
              << "  E=" << c.energy
              << "  survived=" << c.survived
              << "  score="    << c.score  << '\n';
}
\end{lstlisting}

\subsubsection{Relationship to the Identity Vector}

The Hourglass Model consumes $\Sigma_i$ and $\Lambda_i$ directly from
each bead's 32-bit identity word at Gate 1, Gate 3, and the H3 scoring
function.  It is therefore the first model in the platform that makes
the Identity–State Decomposition \emph{operationally load-bearing}
rather than merely descriptive.

% ============================================================================
\section{The Lookglass Bidirectional Feedback Model}
\label{sec:lookglass}
% ============================================================================

\subsection{Concept}

The 6+9 Seed-Bead Stepper is a \emph{forward} pass: atomistic (seed)
information propagates toward the coarse-grained (bead) layer.  This
is the natural direction of coarse-graining.

The Lookglass Model introduces the reverse: bead-layer observables
are reflected back into the atomistic layer as soft corrections.
The name is deliberately evocative — just as a lookglass
(Lewis Carroll, \textit{Through the Looking-Glass}) shows a mirror-reversed
world, the Lookglass pass shows the atomistic layer a reversed image of
its own coarse-grained representation.

\begin{tcolorbox}[
  colback=lgmint, colframe=green!40!black,
  title={Lookglass Architecture},
  fonttitle=\bfseries\small
]
\begin{center}
\begin{tabular}{ll}
\textbf{Forward}  & [S1--S6] $\to$ [B1--B9] $\to$ \texttt{SeedBeadStepRecord}\\[4pt]
\textbf{Backward} & [L1] pos correction $\leftarrow$ bead COM \\
                  & [L2] damping modulation $\leftarrow$ $\bar\eta$ \\
                  & [L3] force scale $\leftarrow$ role weights \\
                  & [L4] dt guard $\leftarrow$ $\bar\Lambda$ \\
                  & [L5] convergence mirror
\end{tabular}
\end{center}
\end{tcolorbox}

\subsection{Backward-Pass Unit Equations}

\subsubsection{L1 — Position Correction}

The CG centroid for bead group $i$ is $\mathbf{r}_i^{\rm bead}$.
The atomistic COM is $\mathbf{r}_i^{\rm seed}$.  The correction:
%
\begin{equation}
  \Delta\mathbf{r}_i^{\rm seed}
  = \alpha_{\rm pos}\,\bigl(\mathbf{r}_i^{\rm bead} - \mathbf{r}_i^{\rm seed}\bigr)
  \label{eq:L1}
\end{equation}
%
$\alpha_{\rm pos} \in [0,1]$ is the position coupling constant.
At $\alpha_{\rm pos}=0$ the backward pass is decoupled; at
$\alpha_{\rm pos}=1$ atomistic positions are snapped to bead centroids.

\subsubsection{L2 — Damping Modulation}

The mean slow state $\bar\eta$ from the BEAD layer inflates the FIRE
damping coefficient:
%
\begin{equation}
  \gamma_{\rm eff} = \gamma_0 \cdot \bigl(1 + \alpha_{\rm damp} \cdot \bar\eta\bigr)
  \label{eq:L2}
\end{equation}
%
High $\bar\eta$ (ordered environment) → larger damping → slower
but more stable relaxation.

\subsubsection{L3 — Force Scale Correction}

For each bead $i$ with dominant role weight $w_{\rm dom}^{(i)}$:
%
\begin{equation}
  F_i^{\rm seed} \;\leftarrow\;
  F_i^{\rm seed} \cdot \bigl[1 + \alpha_{\rm force}(w_{\rm dom}^{(i)} - 1)\bigr]
  \label{eq:L3}
\end{equation}
%
Chemically specific beads ($w_{\rm dom} > 1$) have their forces
amplified; generic Mixed beads ($w_{\rm dom} = 1$) are unchanged.

\subsubsection{L4 — FIRE Step-Size Guard}

Mean stability class $\bar\Lambda \in [0,3]$ modulates the FIRE
timestep ceiling:
%
\begin{equation}
  \Delta t_{\rm max} = \Delta t_0 \cdot
  \operatorname{clamp}\!\left(
    1 - \alpha_{\Delta t}\,\frac{3-\bar\Lambda}{3},\;
    0.1,\; 1.0
  \right)
  \label{eq:L4}
\end{equation}
%
Unstable systems ($\bar\Lambda \approx 0$) reduce the timestep to at
most $10\%$ of $\Delta t_0$; bulk-lattice systems ($\bar\Lambda = 3$)
run at the full timestep.

\subsubsection{L5 — Convergence Mirror}

The backward pass declares steady state when both:
%
\begin{equation}
  \max_i |\Delta\mathbf{r}_i^{\rm feedback}| < r_{\rm tol}
  \quad\text{AND}\quad
  |\bar\eta - \bar\eta_{\rm prev}| < \eta_{\rm tol}
  \label{eq:L5}
\end{equation}
%
Once L5 is satisfied the backward pass is \emph{frozen}: subsequent
ticks skip the backward computation, recovering the pure 6+9 forward
cost.

\subsection{Lookglass Steady State}

\begin{definition}[Lookglass Steady State]
  The combined forward–backward system is in \emph{Lookglass steady
  state} when the 6+9 Seed-Bead Stepper satisfies its own steady-state
  condition (§5.1) \emph{and} the Lookglass backward pass satisfies
  equation~(\ref{eq:L5}).
\end{definition}

The joint steady state is strictly stronger than either condition alone:
it requires that the coarse-grained representation has converged
\emph{and} that its reflection no longer perturbs the atomistic layer.

\subsection{Parameters}

\begin{center}
\begin{tabular}{llrl}
\toprule
Parameter & Symbol & Default & Meaning \\
\midrule
$\alpha_{\rm pos}$ & — & 0.05 & Position coupling (L1) \\
$\alpha_{\rm damp}$ & — & 0.30 & $\bar\eta$ damping inflation (L2) \\
$\gamma_0$ & — & 0.10 & Base FIRE damping (L2) \\
$\alpha_{\rm force}$ & — & 0.20 & Role-weight force scale (L3) \\
$\alpha_{\Delta t}$ & — & 0.50 & Stability-class dt guard (L4) \\
$\Delta t_0$ & — & 5\,fs & Base FIRE timestep (L4) \\
$r_{\rm tol}$ & — & 0.001\,\AA & Feedback displacement tol (L5) \\
$\eta_{\rm tol}$ & — & $10^{-4}$ & Slow-state convergence tol (L5) \\
\bottomrule
\end{tabular}
\end{center}

\subsection{Implementation}

The full implementation is in
\file{coarse\_grain/models/lookglass\_model.hpp}.

\begin{lstlisting}
LookglassParams lg_params;
lg_params.alpha_pos   = 0.05;
lg_params.alpha_force = 0.2;

LookglassRunRecord run_record;
run_record.record_history = true;

double eta_prev = 0.0;

for (uint64_t step = 0; step < max_steps; ++step) {
    // Forward: 6+9 Seed-Bead step
    SeedBeadStepRecord fwd = step_seed_bead(sys, sb_params);

    // Backward: Lookglass reflection
    LookglassStepRecord bwd =
        step_lookglass(sys, fwd, lg_params, eta_prev);
    run_record.ingest(bwd);

    eta_prev = fwd.avg_eta;

    if (fwd.steady_state && bwd.lookglass_steady) break;
}

std::cout << "Lookglass steady at step "
          << run_record.steps_to_steady << '\n';
std::cout << "Peak feedback |dr| = "
          << run_record.peak_pos_correction << " A\n";
\end{lstlisting}

% ============================================================================
\section{Relationships Between the Three Models}
\label{sec:relationships}
% ============================================================================

The three models introduced in this document form a coherent stack:

\begin{center}
\begin{tabular}{lll}
\toprule
Model & Consumes & Produces \\
\midrule
32-Bit Word & $\boldI_i$ fields & \texttt{Identity32} word \\
Hourglass & Bead population + $\Sigma_i$, $\Lambda_i$ & Ranked candidate set \\
Lookglass & Forward step record + $\bar\eta$, $\bar\Lambda$ & Corrected system state \\
\bottomrule
\end{tabular}
\end{center}

\subsection{Interaction Pattern}

A production simulation loop using all three models proceeds as follows:

\begin{enumerate}[label=\textbf{\arabic*.}]
  \item \textbf{Pack identity words.}
        At initialisation (and after any structural event), pack each
        bead's $\boldI_i$ into an \texttt{Identity32} word for fast
        lookup during force evaluation.

  \item \textbf{Forward pass.}
        Run the 6+9 Seed-Bead Stepper.  Channel weights use
        $\Sigma_i$ directly from the bead structs.

  \item \textbf{Backward pass.}
        Run \texttt{step\_lookglass()} using the forward record.
        This corrects bead positions, adjusts FIRE damping and
        timestep, and checks L5.

  \item \textbf{Hourglass probe (periodic).}
        Every $N_{\rm hg}$ steps, run \texttt{run\_hourglass()} to
        probe the local configuration space.  If the best candidate
        scores substantially higher than the current state, accept
        the move.

  \item \textbf{Check joint steady state.}
        Halt when both the Seed-Bead Stepper and the Lookglass
        backward pass have converged.
\end{enumerate}

\subsection{Anti-Black-Box Guarantee}

All three models honour the project-wide anti-black-box contract:

\begin{itemize}
  \item \texttt{Identity32::to\_string()} decodes every bit field by name.
  \item \texttt{CandidateRecord} exposes gate verdicts for every candidate.
  \item \texttt{LookglassStepRecord} exposes every backward-unit output.
\end{itemize}

Every decision made by these models can be reproduced, inspected, and
reported without any hidden state.

% ============================================================================
\section*{Conclusion}
% ============================================================================

This document introduced three foundational models that build directly
on the §0 Identity–State Decomposition framework:

\begin{enumerate}
  \item The \textbf{32-Bit Identity Word} compresses $\boldI_i$ into a
        single register-width integer, enabling branch-free field read-outs
        with no information loss.

  \item The \textbf{Hourglass Model} operationalises $\Sigma_i$ and
        $\Lambda_i$ as load-bearing filter criteria, transforming the
        identity framework from a labelling system into a genuine
        constraint engine for structural convergence.

  \item The \textbf{Lookglass Model} closes the loop between atomistic
        and coarse-grained layers, making the 6+9 Stepper bidirectional
        and enabling joint steady-state convergence.

\end{enumerate}

Together they move VSEPR-SIM from a forward-only simulation pipeline
toward a fully coupled, inspectable, and deterministic research platform.

\end{document}
